\chapter{Шпаргалка}

\section{Verilog на одной странице}

\begin{vcode}
module name #(parameter integer N = 8) (   // параметр
    input  wire         clk,
    input  wire [N-1:0] a,
    output wire         y,
    output reg  [N-1:0] q
);
    localparam [1:0] S0 = 2'd0;            // константа

    wire t = &a;                           // свёртка И
    assign y = t ? 1'b1 : 1'b0;            // мультиплексор

    always @(*) begin                      // комбинационная логика: =
        case (a[1:0])
            2'd0:    ...;
            default: ...;
        endcase
    end

    always @(posedge clk) begin            // триггеры: <=
        if (a == 0) q <= 0;
        else        q <= q + 1'b1;
    end

    other_module u1 (.in(a), .out(...));   // экземпляр модуля
endmodule
\end{vcode}

\section{Типичные ошибки}

\begin{longtable}{p{0.42\textwidth} p{0.52\textwidth}}
\toprule
\textbf{Симптом} & \textbf{Вероятная причина} \\
\midrule
\endhead
Светодиоды горят «наоборот» & Забыли, что LED горят от 0: нужна инверсия \\
Ничего не работает, схема «пустая» & Выход модуля ни к чему не подключён, и синтезатор удалил всю логику; неверные имена в \code{.cst} \\
Счётчик считает одно нажатие много раз & Нет антидребезга и/или детектора фронта \\
Предупреждение \emph{latch inferred} & В \code{always @(*)} не хватает \code{else} или \code{default} \\
Два регистра не обмениваются значениями & В \code{always @(posedge clk)} использовано \code{=} вместо \code{<=} \\
Число не помещается, счётчик не доходит до нужного значения & Мала разрядность регистра: проверьте $2^n > N$ \\
\emph{Multiple drivers} & Одному сигналу присваивают значение в двух разных \code{always} или \code{assign} \\
После выключения питания схема пропадает & Загрузили в SRAM; для постоянной работы загрузите во Flash \\
\bottomrule
\end{longtable}

\section{Выводы Tang Nano 20K}

\begin{table}[H]
\centering
\begin{tabular}{l c l}
\toprule
\textbf{Сигнал} & \textbf{Вывод} & \textbf{Назначение} \\
\midrule
\code{clk} & 4 & 27 МГц \\
\code{btn[0]}, \code{btn[1]} & 88, 87 & кнопки S1, S2 \\
\code{led[0..5]} & 15--20 & светодиоды (активный 0) \\
\code{uart\_tx}, \code{uart\_rx} & 69, 70 & UART через USB \\
\code{ws2812} & 79 & RGB-светодиод \\
\bottomrule
\end{tabular}
\end{table}

\section{Полезные ссылки}

\begin{itemize}
  \item Документация Sipeed: \url{https://wiki.sipeed.com/hardware/en/tang/tang-nano-20k/nano-20k.html}
  \item Примеры Sipeed: \url{https://github.com/sipeed/TangNano-20K-example}
  \item Среда Gowin EDA: \url{https://www.gowinsemi.com}
  \item Открытые инструменты (OSS CAD Suite): \url{https://github.com/YosysHQ/oss-cad-suite-build}
  \item Проект Apicula: \url{https://github.com/YosysHQ/apicula}
  \item Интерактивный учебник по Verilog HDLBits: \url{https://hdlbits.01xz.net}
\end{itemize}
